\chapter{Ответы и подсказки к упражнениям}

\subsection*{Глава 1}
\textbf{1.1.} $\Delta G = \Delta H\Delta T/T_e$: 100 и 1000~Дж/моль. $RT$ при 990 и 900~К --- 8.2 и 7.5~кДж/моль, то есть
движущая сила меньше $RT$ в 82 и 7.5 раза.
\textbf{1.2.} $\exp\bigl[\frac{Q}{R}\bigl(\frac1{573}-\frac1{673}\bigr)\bigr] = \exp(3.74)\approx42$.
\textbf{1.3.} Вдоль линии равновесия $\dd G_\alpha = \dd G_\beta$, то есть $V_\alpha\dd p - S_\alpha\dd T = V_\beta\dd p -
S_\beta\dd T$, откуда $\dd p/\dd T = \Delta S/\Delta V = \Delta H/(T\Delta V)$. Для напряжения роль $-p\,\dd V$ играет работа
$V\sigma:\dd\eps$, и $\dd T/\dd\sigma = T\eps^*V/\Delta H$ --- формула главы~6.

\subsection*{Глава 2}
\textbf{2.1.} $\mu_B = G - (X-1)\,\partial G/\partial X$ даёт $G_B + \Omega(1-X)^2 + RT\ln X$.
\textbf{2.2.} $Q = 2.24\cdot10^3\cdot\ln10\cdot R \approx 42.9$~кДж/моль. $T = 2240/(5.56-\lg C)$: 100~ppm --- 356\degC,
300~ppm --- 453\degC, 600~ppm --- 532\degC (почти эвтектоид).
\textbf{2.3.} $\Delta g = R\ln10\,(560 - 0.48\,T)$~Дж/моль: 6.4~кДж/моль при 200\degC, 5.5 при 300\degC, 4.5 при
400\degC. Падает с температурой: при высокой $T$ пластичность и возврат легче, и матрица «берёт» с гидрида меньше.
\textbf{2.4.} $C_\TSSD(400\degC)\approx171$~ppm растворено; около 129~ppm остаётся в гидридах. Выпадение начнётся на TSSP
при 171~ppm: $T = 1680/(5.08-\lg171)\approx 317\degC$.

\subsection*{Глава 3}
\textbf{3.1.} $D(350\degC)\approx1.5\cdot10^{-10}$~м$^2$/с. $t\sim L^2/D$: для 0.65~мм около 50~мин, для 50~мкм около 17~с.
\textbf{3.2.} $\mathbf J = 0$ при $\nabla\ln c = (\bar V/RT)\nabla\sigma_h$. При 600~МПа и 300\degC: $e^{0.21}\approx1.24$.
\textbf{3.3.} $g = 1 - 4r\sin^2(k\Delta x/2)$; $|g|\le1$ при $r\le1/2$.

\subsection*{Глава 4}
\textbf{4.1.} $\operatorname{tg}\varphi = 1.1547\cdot1.593/7 = 0.263$, $\varphi\approx14.7^\circ$.
\textbf{4.2.} $1.17^{1/3}-1\approx0.054$; дислокация примерно через каждые 18 плоскостей.
\textbf{4.3.} $\mu b^2/2\approx1.7\cdot10^{-9}$~Дж/м, то есть около 3.5~эВ на межатомное расстояние.

\subsection*{Глава 5}
\textbf{5.1.} $\sigma^{\text{вн}} = 3K(s-1)\eps\,\mathbf I$, $s = \frac{1+\nu}{3(1-\nu)}$; $W = \frac92K(1-s)\eps^2V = \frac{E}{1-\nu}\eps^2V$.
\textbf{5.2.} В пластинке $\eps^{\text{полн}}_{tt} = \eps^{\text{полн}}_{ss} = 0$, $\sigma_{nn} = \sigma_{nt} = \sigma_{ns} = 0$:
плоское напряжённое состояние с упругой деформацией $-\eps^*$ в плоскости; $w = \frac12\sigma:\eps_{\text{упр}}$.
\textbf{5.3.} $\eps_{tt} = 0.0476$; $w\approx297$~МДж/м$^3$, около 3.0~кДж/моль~H.
\textbf{5.4.} Ось $c$ по окружности: $(120\cdot0.072 + 60\cdot0.046)\cdot10 \approx 114$~Дж/моль. Ось $c$ по радиусу:
$(120\cdot0.046 + 60\cdot0.046)\cdot10\approx 82$~Дж/моль.

\subsection*{Глава 6}
\textbf{6.1.} $r^* = 2\gamma/\Delta g$, $\Delta G^* = 16\pi\gamma^3/(3\Delta g^2)$; при $\Delta g\to0.8\Delta g$ барьер растёт в
$1/0.64\approx1.56$ раза.
\textbf{6.2.} Объём линзы $V = \frac43\pi R^3 S(\theta)$, площадь межфазной поверхности и заменённой границы дают тот же
множитель в барьере.
\textbf{6.3.} $\delta T\approx0.013\cdot150\approx2$~К по упругой оценке, $\approx12$~К по измеренному сдвигу; полоса гистерезиса ---
около 80~К.
\textbf{6.4.} $f_{\text{расш}} = N\frac43\pi v^3t^3$, $n = 3$.

\subsection*{Глава 7}
\textbf{7.1.} Мизес: $\sigma_\theta = 2\sigma_y/\sqrt3\approx1.155\sigma_y$; Треска ($\sigma_r\approx0$): $\sigma_\theta = \sigma_y$; разница 15\,\%.
\textbf{7.2.} $\sigma_{rr} = -4\mu B/r^3$, $\sigma_{\theta\theta} = 2\mu B/r^3$; текучесть на $r = c$: $6\mu B/c^3 = \sigma_y$;
несжимаемость: $a^2u(a) = c^2u(c) = B$, $u(a) = \delta a$.
\textbf{7.3.} $f_r = \cos^2\chi$, $f_\theta = \sin^2\chi$; равны при $\chi = 45^\circ$.
